\section{Eksport Model Terbaik}

Berdasarkan hasil perbandingan evaluasi kinerja sebelumnya, model \textbf{\textit{XGBoost} versi tanpa \textit{tuning}} terbukti menghasilkan tingkat akurasi maksimal dan metrik \textit{Recall} untuk kelas minoritas yang sempurna secara efisien dengan waktu eksekusi yang paling cepat. Oleh karena itu, model \textit{XGBoost} yang diimplementasikan pada sistem. Proses ekstraksi dan penyimpanan (\textit{export}) model dilakukan menggunakan pustaka \texttt{joblib} pada bahasa pemrograman Python. Pada tahap ini, model \textit{XGBoost} digabungkan menjadi satu kesatuan bersama komponen pendukung utamanya. Komponen pendukung tersebut meliputi \textit{Label Encoder} yang berfungsi untuk melakukan transformasi dari nilai kategorikal teks menjadi format angka, serta aturan pendefinisian tiga fitur terpilih. Penggabungan komponen ini bertujuan untuk memastikan bahwa setiap data masukan (\textit{input}) baru yang dikirimkan melalui aplikasi dapat diproses dan diseragamkan formatnya secara otomatis agar sama persis dengan struktur data yang digunakan pada fase pelatihan model, sehingga mencegah terjadinya \textit{error} akibat ketidaksesuaian tipe data. Kesatuan objek tersebut kemudian disimpan ke dalam format berkas \textit{pickle} dengan nama \texttt{model\_xgboost.pkl}.

Berkas model yang telah diekstrak tersebut selanjutnya diintegrasikan secara langsung ke dalam aplikasi Data Capture. Melalui integrasi ini, aplikasi mampu memberikan klasifikasi secara langsung (\textit{real-time}) terhadap setiap data penyewa baru yang dimasukkan. Sistem akan menampilkan peringatan berupa pesan ``Penyewa Berpotensi Melakukan Pelanggaran!'' apabila seorang penyewa terklasifikasi ke dalam kelas 1, yang mengindikasikan penyewa berpotensi melakukan pelanggaran. Sebaliknya, apabila profil data penyewa yang diinput tergolong ke dalam kelas 0 (terindikasi normal dan aman), sistem tidak akan memunculkan peringatan apa pun pada layar antarmuka pengguna.